\begin{tresc}
Szachownica o wymiarach 3×3 ma cztery pola białe i pięć czarnych. W każdym ruchu przemalowujemy dowolnie wybrane trzy pola (zawsze przemalowując pola czarne na biało, a białe na czarno). Jaka jest najmniejsza liczba ruchów, przy pomocy których można doprowadzić do sytuacji, gdy wszystkie pola będą białe? Każde pole może być przemalowywane wielokrotnie.
\end{tresc}
\begin{wariant}{A} 2 \end{wariant}
\begin{wariant}{B} 3 \end{wariant}
\begin{wariant}{C} 4 \end{wariant}
\begin{wariant}{D} 5 \end{wariant}
\begin{wariant}{E} nie da się tego wykonać \end{wariant}
